\chapter{Инвариантное ядро}

Рабочая гипотеза единого источника требует указать, что именно сохраняется
при смене языка описания. Под инвариантным ядром $\mathcal I(\mathfrak S)$
будем понимать совокупность свойств, которые переживают допустимые
геометрические, спектральные и корреляционные редукции.

Ядро не является перечнем всех формул проекта. В него входят только те
структуры, без которых теряются одновременно:
\begin{itemize}
 \item идентичность предполагаемого источника $\mathfrak S$;
 \item различимость способов его чтения;
 \item управляемый переход к наблюдаемым;
 \item возможность сравнивать конкурирующие модели.
\end{itemize}

Минимальный рабочий каркас записывается как
\begin{equation}
 \mathfrak K_{\min}=
 \bigl(\mathfrak S,\mathcal I(\mathfrak S),
 \mathcal R_{\mathrm{type}},\mathcal Q_{\mathrm{sel}}\bigr),
\end{equation}
где $\mathcal R_{\mathrm{type}}$ задаёт типы редукции, а
$\mathcal Q_{\mathrm{sel}}$ --- принцип выбора между моделями.

\section{Проверка исключением}

Принадлежность элемента ядру устанавливается не декларацией, а проверкой:
\begin{enumerate}
 \item элемент временно исключается из аппарата;
 \item проверяется сохранение идентичности $\mathfrak S$;
 \item проверяется различимость режимов чтения;
 \item проверяется возможность реконструкции наблюдаемых;
 \item проверяется сохранение нетривиального выбора между моделями.
\end{enumerate}

Если все функции сохраняются, элемент считается выводимым или вспомогательным.
Если хотя бы одна несущая функция исчезает, элемент остаётся кандидатом в
фундаментальную часть ядра. При недостатке данных фиксируется неопределённый
статус, а не преждевременный положительный вывод.

\chapter{Секторное разложение чтения}

Для первичной классификации вводятся четыре вклада:
\begin{equation}
 W_{\mathrm{car}},\qquad W_{\mathrm{phase}},\qquad
 W_{\mathrm{int}},\qquad W_{\mathrm{mix}}.
\end{equation}

Они обозначают соответственно каркасный, фазовый, внутренний и смешанный
способы формирования наблюдаемого. После выбора положительной нормировки
вводятся веса
\begin{equation}
 w_a=\frac{W_a}{
 W_{\mathrm{car}}+W_{\mathrm{phase}}+W_{\mathrm{int}}+W_{\mathrm{mix}}},
 \qquad
 \sum_a w_a=1.
\end{equation}

Эти веса не считаются свободными физическими параметрами. Они допустимы лишь
как результаты одной заранее заданной процедуры чтения. Если для разных
наблюдаемых требуется независимо выбирать нормировку или правило разложения,
секторная атрибуция считается послепроверочной и не проходит отбор.

\chapter{Типы структурных режимов}

По преобладающему вкладу различаются следующие режимы.

\section{Каркасный режим}

Наблюдаемое определяется главным образом геометрией, топологией или ведущим
спектральным законом носителя. Его значение должно быть устойчиво к слабым
изменениям фазового и внутреннего секторов.

\section{Фазовый режим}

Наблюдаемое чувствительно к выбору глобальной ветви, голономии, спиновой
структуры или ориентированной фазовой информации. Локальное изменение без
смены глобального класса не должно воспроизводить тот же эффект произвольно.

\section{Внутренний режим}

Наблюдаемое определяется внутренним операторным содержанием, кратностями,
матрицами связи или иерархией состояний. Его нельзя выводить только из объёма
или топологии внешнего носителя.

\section{Смешанный режим}

Наблюдаемое существенно зависит не менее чем от двух секторов, причём вклад
не раскладывается в сумму независимо настраиваемых частей. Для такого режима
необходим общий оператор или единый функционал, порождающий смешанный член.

\section{Переходный режим}

При допустимой деформации наблюдаемое может менять преобладающий тип. Такой
переход должен описываться непрерывной траекторией или явно типизированным
скачком, а не переименованием режима после получения результата.

\chapter{Карта структурных режимов}

Для качественной карты вводятся две предварительные координаты. Величина
$\Xi$ измеряет удаление от чисто каркасного режима, а $\Upsilon$ различает
фазовую и внутреннюю ориентацию некаркасной части:
\begin{equation}
 \mathfrak R(O)=\bigl(\Xi(O),\Upsilon(O)\bigr).
\end{equation}

При $\Xi\ll1$ наблюдаемое относится к каркасной области. При больших
$\Xi$ знак $\Upsilon$ различает преимущественно фазовую и преимущественно
внутреннюю области, а $\Upsilon\approx0$ указывает на смешанный режим.

Координаты $\Xi$ и $\Upsilon$ имеют пока методический статус. До вывода их
единственной нормировки они не являются новыми физическими наблюдаемыми и не
могут использоваться для численного предсказания.

\chapter{Правило секторной атрибуции}

Отнесение наблюдаемого $O$ к режиму проводится в пять шагов:
\begin{enumerate}
 \item указать оператор или функционал, порождающий $O$;
 \item определить допустимые деформации каждого сектора;
 \item вычислить чувствительность $O$ к этим деформациям;
 \item проверить устойчивость классификации при смене представления;
 \item зафиксировать условие, при котором классификация признаётся неверной.
\end{enumerate}

Одного словесного сходства с геометрией, фазой или внутренней структурой
недостаточно. Атрибуция должна следовать из поведения наблюдаемого под
контролируемыми изменениями модели.

\chapter{Ограничение карты}

Карта режимов является средством организации программы, а не доказанной
геометрией пространства состояний. Её размерность, метрика и границы областей
должны быть выведены из строгой модели. До этого момента карта используется
только для постановки различающих вопросов и поиска скрытых зависимостей.

Следующий сегмент должен определить процедуру реконструкции наблюдаемых и
правила сравнения конкурирующих моделей без подгонки секторных коэффициентов.